Puisqu'on interroge une personne au hasard, on est en situation d'équiprobabilité, et les proportions sont donc assimilables à des probabilités.

	\begin{enumerate}
		\item Il y a au total 100 personnes, dont 35 ne vivant pas en couple et n'ayant pas d'enfant,

 		donc \mbox{$P\big(\overline{C}\cap\overline{E}\big) = \dfrac{35}{100} = \np{0.35}$.}

		\item \mbox{$P\big(C\cap\overline{E}\big) = \dfrac{12}{100} = \np{0.12}$} : \quad la probabilité de choisir au hasard une personne vivant en couple et n'ayant pas d'enfant est \np{0.12}.

		\item On cherche \mbox{$P_{E}\big(\overline{C}\big)$}.

		Par définition des probabilités conditionnelles, \mbox{$P_{E}\big(\overline{C}\big) = \dfrac{P\big(E\cap\overline{C}\big)}{P\big(E\big)}$}.

		53 personnes sur les 100 étudiées ont au moins un enfant, donc \mbox{$P\big(E\big) = \np{0.53}$}.

		Parmi les personnes interrogées, 35 ne vivent pas en couple mais ont au moins un enfant, donc \mbox{$P\big(E\cap\overline{C}\big) = \np{0.35}$}

		La probabilité que la personne ne vive pas en couple sachant qu'elle a au moins un enfant est donc \mbox{$P_{E}\big(\overline{C}\big) = \dfrac{35}{53}$.}

		NB : on peut retrouver plus simplement ce résultat en observant que, parmi les 53 personnes ayant au moins un enfant, 35 ne vivent pas en couple.

		\item \begin{enumerate}
			\item Un logement peut donc avoir 1, 2, 3, 4 ou 5 pièces, et les probabilités sont les pourcentages respectifs.

			On a donc la loi de probabilité de $X$ :

			\begin{center}
			\begin{tblr}{colspec={X[l,2cm]*{5}{X[c,1.5cm]}},hlines, vlines}
				$x_i$        & $1$  & $2$ & $3$ & $4$  & $5$  \\
				$P(X = x_i)$ & 0,07 & 0,2 & 0,4 & 0,25 & 0,08\\
			\end{tblr}
			\end{center}

			\item \mbox{$P(X\geqslant 2)$} est l'évènement contraire de \mbox{$P(X = 1)$}, donc \quad \mbox{$P(X\geqslant 2) = 1-\np{0.07} = \np{0.93}$.}

			cela signifie que la probabilité qu'un logement ait un nombre de pièces supérieur ou égal à 2 est \np{0.93}.

			\item L'espérance de $X$ est égale à \quad $\ds \mathrm{E}(X) = \sum_{i = 1}^5 P(X = x_i)\times x_i$.

			$\ds \mathrm{E}(X) = 0,07 \times 1 + 0,2 \times 2 + 0,4 \times 3 + 0,25 \times 4 + 0,08 \times 5 = 0,07 + 0,4 + 1,2 + 1 + 0,4 = 0,47 + 2,2 + 0,4 = 0,47 + 2,6 = 3,07$.

			donc \mbox{$E(X) = \np{3.07}$}.

			Cela signifie que les logements étudiés ont en moyenne 3,07 pièces.

			\emph{Remarque :} Ici, le 3,07 s'interprète comme une moyenne, il n'y a donc pas de raison d'arrondir le résultat à un nombre entier.
		\end{enumerate}
	\end{enumerate}
